\documentclass{article}
\usepackage[preprint]{neurips_2026}
\usepackage{booktabs}
\usepackage{graphicx}
\usepackage{hyperref}
\usepackage{amsmath}

\title{How Well Can RAG Systems Write Literature Reviews?\\A Benchmark for Citation-Grounded Scientific Synthesis}

\author{
  Author One \\
  Affiliation \\
  \texttt{author@institution.edu}
}

\begin{document}
\maketitle

\begin{abstract}
Retrieval-Augmented Generation (RAG) promises grounded, factual text generation
by conditioning on retrieved documents. But can RAG systems produce accurate,
citation-grounded scientific literature reviews comparable to expert-written
surveys? We introduce \textbf{LitSynthBench}, a benchmark of 15 ML/NLP topics
with expert-written survey papers, candidate corpora of 200+ papers per topic,
and human annotations for error classification. We evaluate three RAG
configurations (dense, sparse, hybrid retrieval) with multiple generator models
against a no-retrieval baseline and expert-written reviews. Our evaluation
framework jointly measures citation accuracy, factual consistency, coverage
completeness, and error patterns---the first to assess all four dimensions for
multi-document scientific synthesis. We find that RAG significantly reduces
hallucinations compared to pure generation, but citation accuracy and coverage
remain substantially below expert levels, with systematic biases toward recent
and highly-cited papers. We propose a four-category error taxonomy (citation
hallucination, attribution error, omission, recency bias) with inter-annotator
agreement $\kappa > 0.6$. Our benchmark and evaluation framework provide a
standardized foundation for measuring and improving RAG-based literature
synthesis.
\end{abstract}

% ============================================================================
\section{Introduction}
% ============================================================================

The use of large language models (LLMs) for scientific writing is rapidly
growing~\cite{gao2024rag_survey}. Retrieval-Augmented Generation (RAG) systems
combine document retrieval with language generation to produce grounded
outputs~\cite{lewis2020rag}. While RAG has shown strong results on
question-answering and summarization tasks, its application to multi-document
scientific literature synthesis---a task requiring accurate citations,
comprehensive coverage, and faithful attribution---remains under-evaluated.

Existing benchmarks evaluate RAG on single-document QA
(MIRAGE~\cite{xiong2024mirage}), faithfulness in
summarization~\cite{faithjudge2025}, or general hallucination
detection~\cite{confabulations2024}. None provide a standardized evaluation for
\emph{citation-grounded literature synthesis}, where the system must survey a
body of work, identify key contributions, and correctly attribute claims to
specific papers.

We address this gap with three contributions:

\begin{enumerate}
    \item \textbf{LitSynthBench}: A benchmark of 15 ML/NLP topics with
    expert-written survey papers, retrieval corpora, and human annotations.
    \item \textbf{A comprehensive evaluation framework} jointly measuring
    citation accuracy, factual consistency, coverage, and error patterns.
    \item \textbf{An error taxonomy} categorizing RAG synthesis failures into
    citation hallucination, attribution error, omission, and recency bias.
\end{enumerate}

% ============================================================================
\section{Related Work}
% ============================================================================

\subsection{Retrieval-Augmented Generation}
% TODO: Expand with actual citations from literature-review.md

\subsection{Hallucination Detection and Faithfulness}
% TODO: FaithJudge, HHEM, RAGTruth

\subsection{Automated Literature Review Generation}
% TODO: arxiv:2411.18583, arxiv:2407.20906, EMNLP 2025

\subsection{Evaluation Benchmarks}
% TODO: MIRAGE, Confabulations, FaithJudge

% ============================================================================
\section{LitSynthBench}
% ============================================================================

\subsection{Topic Selection}
We select 15 sub-topics from ML/NLP with published expert survey papers from
top venues (ACL, EMNLP, NeurIPS, ICML) between 2022--2025.

\subsection{Corpus Construction}
For each topic, we construct a retrieval corpus using the Semantic Scholar
API~\cite{semantic_scholar}: 200--500 candidate papers retrieved by topic query,
plus all papers cited by the expert survey.

\subsection{Annotation Protocol}
Two annotators classify errors in RAG-generated reviews into four categories.
Target inter-annotator agreement: Cohen's $\kappa > 0.6$.

% ============================================================================
\section{Experimental Setup}
% ============================================================================

\subsection{Conditions}
% Table from experiment-plan.md
\begin{table}[h]
\centering
\caption{Experimental conditions.}
\begin{tabular}{llll}
\toprule
ID & System & Retrieval & Generator \\
\midrule
C1 & RAG + Dense & Contriever/BGE & LLaMA-3.1-8B \\
C2 & RAG + Sparse & BM25 & LLaMA-3.1-8B \\
C3 & RAG + Hybrid & BM25 + Dense rerank & LLaMA-3.1-8B \\
C4 & No Retrieval & None & LLaMA-3.1-8B \\
C5 & Expert & N/A & Human \\
\bottomrule
\end{tabular}
\end{table}

\subsection{Metrics}
Citation precision and recall, faithfulness score, hallucination rate, key
paper recall (KPR), ROUGE, BERTScore, recency and popularity correlation.

% ============================================================================
\section{Results}
% ============================================================================

\textit{[PLACEHOLDER: Results pending experiment execution.]}

\subsection{Citation Accuracy (H1)}
\subsection{Factual Consistency (H2)}
\subsection{Coverage Analysis (H3)}
\subsection{Error Taxonomy (H4)}
\subsection{Ablation Studies}

% ============================================================================
\section{Discussion}
% ============================================================================

\textit{[PLACEHOLDER]}

% ============================================================================
\section{Conclusion}
% ============================================================================

\textit{[PLACEHOLDER]}

\bibliography{references}
\bibliographystyle{plainnat}

\end{document}
